\documentclass[aps,preprint]{revtex4}%
\usepackage{amsfonts}%
\usepackage{amsmath}%
\setcounter{MaxMatrixCols}{30}%
\usepackage{amssymb}%
\usepackage{graphicx}
%TCIDATA{OutputFilter=latex2.dll}
%TCIDATA{Version=5.50.0.2953}
%TCIDATA{CSTFile=revtex4.cst}
%TCIDATA{Created=Friday, June 20, 2008 10:01:44}
%TCIDATA{LastRevised=Wednesday, June 25, 2008 09:43:38}
%TCIDATA{<META NAME="GraphicsSave" CONTENT="32">}
%TCIDATA{<META NAME="SaveForMode" CONTENT="1">}
%TCIDATA{BibliographyScheme=Manual}
%TCIDATA{<META NAME="DocumentShell" CONTENT="Articles\SW\REVTeX 4">}
%BeginMSIPreambleData
\providecommand{\U}[1]{\protect\rule{.1in}{.1in}}
%EndMSIPreambleData
\newtheorem{theorem}{Theorem}
\newtheorem{acknowledgement}[theorem]{Acknowledgement}
\newtheorem{algorithm}[theorem]{Algorithm}
\newtheorem{axiom}[theorem]{Axiom}
\newtheorem{claim}[theorem]{Claim}
\newtheorem{conclusion}[theorem]{Conclusion}
\newtheorem{condition}[theorem]{Condition}
\newtheorem{conjecture}[theorem]{Conjecture}
\newtheorem{corollary}[theorem]{Corollary}
\newtheorem{criterion}[theorem]{Criterion}
\newtheorem{definition}[theorem]{Definition}
\newtheorem{example}[theorem]{Example}
\newtheorem{exercise}[theorem]{Exercise}
\newtheorem{lemma}[theorem]{Lemma}
\newtheorem{notation}[theorem]{Notation}
\newtheorem{problem}[theorem]{Problem}
\newtheorem{proposition}[theorem]{Proposition}
\newtheorem{remark}[theorem]{Remark}
\newtheorem{solution}[theorem]{Solution}
\newtheorem{summary}[theorem]{Summary}
\newenvironment{proof}[1][Proof]{\noindent\textbf{#1.} }{\ \rule{0.5em}{0.5em}}
\begin{document}
\preprint{HEP/123-qed}
\title[Short title for running header]{Long title that appears on the title page}
\author{First Author}
\affiliation{My Institution}
\author{Second Author}
\affiliation{My Institution}
\author{Third Author}
\affiliation{Other Institution}
\keywords{one two three}
\pacs{PACS number}

\begin{abstract}
Shell document for REV\TeX{} 4.

\end{abstract}
\volumeyear{year}
\volumenumber{number}
\issuenumber{number}
\eid{identifier}
\date[Date text]{date}
\received[Received text]{date}

\revised[Revised text]{date}

\accepted[Accepted text]{date}

\published[Published text]{date}

\startpage{101}
\endpage{102}
\maketitle
\tableofcontents
%

\[
E_{H}\left(  H\right)  =Tr\left\{  HD\right\}
\]%
\[
E_{GS}=E_{H}\left(  D_{\#}\right)  =\min_{D\in S_{N}}Tr\left\{  HD\right\}
\]%
\[
S_{N}=\left\{  \left.  D\right\vert D\in\mathcal{B}_{1}\left(  \mathcal{H}%
\right)  \right\}
\]
$S_{N}$ is closed in the topolgy. If $H$ is bounded then $E_{H}\in
\mathcal{B}\left(  \mathcal{H}\right)  $ and is continuous and bounded, if $H$
is not bounded but densely defined then $E_{H}$ is not continuous on its dense
domain and not bounded.
\[
\mathcal{B}\left(  \mathcal{H}\right)  ^{\ast}\supset\mathcal{B}_{1}\left(
\mathcal{H}\right)
\]
As
\[
\lim\mathcal{B}_{1}\left(  \mathcal{H}\right)  \overset{\text{Operator Norm}%
}{\longrightarrow}\mathcal{B}\left(  \mathcal{H}\right)
\]
but
\[
\lim\mathcal{B}_{1}\left(  \mathcal{H}\right)  \overset{\text{weak *}%
}{\longrightarrow}\mathcal{B}\left(  \mathcal{H}\right)  ^{\ast}%
\]%
\begin{align*}
\mathcal{B}\left(  \mathcal{H}\right)  ^{\ast}  & \supset\mathcal{B}%
_{1}\left(  \mathcal{H}\right)  \Leftrightarrow\mathcal{B}_{1}\left(
\mathcal{H}\right)  ^{\ast\ast}\supset\mathcal{B}_{1}\left(  \mathcal{H}%
\right)  \\
& \Rightarrow\mathcal{B}\left(  \mathcal{H}\right)  ^{\ast\ast}\subset
\mathcal{B}_{1}\left(  \mathcal{H}\right)  ^{\ast}\\
& =\mathcal{B}\left(  \mathcal{H}\right)  ^{\ast\ast}\subset\mathcal{B}\left(
\mathcal{H}\right)
\end{align*}%
\[
\mathcal{B}\left(  \mathcal{H}\right)  ^{\ast\ast}\supset\mathcal{B}\left(
\mathcal{H}\right)
\]%
\[
\Rightarrow\mathcal{B}\left(  \mathcal{H}\right)  ^{\ast\ast}=\mathcal{B}%
\left(  \mathcal{H}\right)
\]%
\[
H=\int_{\sigma\left(  H\right)  }EdP_{H}\left(  E\right)  =\int_{\sigma\left(
H\right)  }E\left\vert \Psi_{H}\left(  E\right)  \right\rangle \left\langle
\Psi_{H}\left(  E\right)  \right\vert dE
\]%
\[
\left\vert \Psi\left(  E\right)  \right\rangle \in\mathcal{D}_{H}^{\ast
}\supset\mathcal{H\supset D}_{H}%
\]
Many different riggings.%
\[
\boldsymbol{r}=\int_{\sigma\left(  Q\right)  }\boldsymbol{r}dP_{Q}\left(
\boldsymbol{r}\right)  =\int_{\sigma\left(  Q\right)  }\boldsymbol{r}%
\left\vert \Psi_{Q}\left(  \boldsymbol{r}\right)  \right\rangle \left\langle
\Psi_{Q}\left(  \boldsymbol{r}\right)  \right\vert d\boldsymbol{r}%
\]%
\[
\left\vert \Psi\left(  \boldsymbol{r}\right)  \right\rangle \in\mathcal{D}%
_{Q}^{\ast}\supset\mathcal{H\supset D}_{Q}%
\]%
\[
\rho=\int_{\sigma\left(  Q\right)  }\rho\left(  \boldsymbol{r}\right)
dP_{Q}\left(  \boldsymbol{r}\right)  =\int_{\sigma\left(  Q\right)  }%
\rho\left(  \boldsymbol{r}\right)  \left\vert \Psi_{Q}\left(  \boldsymbol{r}%
\right)  \right\rangle \left\langle \Psi_{Q}\left(  \boldsymbol{r}\right)
\right\vert d\boldsymbol{r}%
\]%
\[
\rho=\int\Phi\left(  \boldsymbol{r}\right)  ^{\dagger}\Phi\left(
\boldsymbol{r}\right)  d\left(  \boldsymbol{r}\right)
\]%
\begin{equation}
\mathcal{B}_{2}\left(  \mathcal{H}^{N}\right)  =\mathcal{B}_{2sa}\left(
\mathcal{H}^{N}\right)  \oplus\mathcal{B}_{2asa}\left(  \mathcal{H}%
^{N}\right)
\end{equation}
%

\[
P_{+}\left(  X\right)  =\frac{1}{2}\left\{  X+X^{\dagger}\right\}  \Rightarrow
P_{+}\left(  X\right)  ^{\dagger}=P_{+}\left(  X\right)
\]
%

\[
P_{-}\left(  X\right)  =\frac{1}{2}\left\{  X-X^{\dagger}\right\}  \Rightarrow
P_{-}\left(  X\right)  ^{\dagger}=-P_{-}\left(  X\right)
\]%
\[
P_{+}\left(  X\right)  +P_{-}\left(  X\right)  =\left(  P_{+}+P_{-}\right)
\left(  X\right)  =X
\]%
\[
L_{1}\left(  \mathbb{R}^{3}\right)  \cap L_{3}\left(  \mathbb{R}^{3}\right)
\supset L_{2}\left(  \mathbb{R}^{3}\right)
\]%
\[
L_{1}\left(  \mathbb{R}^{3}\right)  \cap L_{3}\left(  \mathbb{R}^{3}\right)
\supset L_{2}\left(  \mathbb{R}^{3}\right)
\]%
\begin{align*}
\left[  L_{1}\left(  \mathbb{R}^{3}\right)  \cap L_{3}\left(  \mathbb{R}%
^{3}\right)  \right]  ^{\ast}  & \subset L_{2}\left(  \mathbb{R}^{3}\right)
^{\ast}\\
L_{\infty}\left(  \mathbb{R}^{3}\right)  \oplus L_{\frac{3}{2}}\left(
\mathbb{R}^{3}\right)    & \subset L_{2}\left(  \mathbb{R}^{3}\right)
\end{align*}%
\[
\int a\left(  \boldsymbol{r}\right)  \left\vert \boldsymbol{r}\right\rangle
\left\langle \boldsymbol{r}\right\vert d\boldsymbol{r=}\int b\left(
\boldsymbol{r}\right)  \left\vert \boldsymbol{r}\right\rangle \left\langle
\boldsymbol{r}\right\vert d\boldsymbol{r}%
\]%
\[
f=\int f\left(  \boldsymbol{r}\right)  \left\vert \boldsymbol{r}\right\rangle
d\boldsymbol{r}%
\]
All $L_{p}$ spaces are included in the space of distributions.


\end{document}